Neural Networks (NNs) are typically trained by first fixing the architecture $A \in \mathcal{A}$, after which the parameters $\theta \in \Theta_A$ are optimized to minimize a training objective. The architectural choice is crucial, as it defines the function class explored during training, 
and is the workhorse of deep learning~\citep{he_deep_2016,vaswani_attention_2017}. Despite this, architecture search is still typically performed manually, requiring significant expertise and wasted compute due to retraining.

Ideally, we would like to jointly optimize over the architecture space $\mathcal{A}$ and its space of parameters $\Theta_A$ by solving
\begin{align}\label{eqn:ideal_obj}
    A^*, \theta^* = \argmin_{A\in \mathcal{A},\; \theta \in  \Theta_A}  \mathcal{L}(f_{A,\theta}) + \mathcal{C}(A)
\end{align}
where $f_{A,\theta}$ denotes the function induced by architecture $A$ with parameters $\theta$, $\mathcal{L}$ is the empirical risk over a data distribution $\mathcal{D}$, and $\mathcal{C}$ is a penalty on the architecture complexity, measured in total parameters, FLOPs or some other notion of complexity. The closest approach to this objective is \emph{Neural Architecture Search} (NAS)~\citep{zoph_neural_2017}\footnote{NAS does not exactly optimize this objective, since it typically does not find the optimal parameters $\theta^*$.}
as a form of hyper-parameter optimization~\citep{egele:tel-04636586}. However, a full NAS loop is often prohibitively expensive, requiring multiple retrainings, and also ignores a key constraint: we frequently start from a pre-trained model that we would like to adapt rather than discard.

This motivates \emph{growing} Neural Network architectures: starting with a small ``seed'' architecture and expanding its capacity during training by applying local architecture transformations $\T$ (also known as network morphisms), such as widening existing layers or adding new ones, and appropriately adapting the existing weights. Concretely, let $f_{A_t, \theta_t}$ denote the model at growth step $t$ and $\Opt_
\theta(\textrm{goal}(\theta), \theta_{\text{init}})$ denote a few steps of, e.g., stochastic gradient descent to optimize $\text{goal}(\theta)$ over $\theta$ starting from $\theta_{\text{init}}$.
We alternate between applying the growth operator $\T$ and the optimization step
\begin{equation}\label{eqn:grow_decomposition}
\begin{aligned}
    &\theta_t' = \Opt_{\theta}(\mathcal{L}(f_{A_t, \theta}), \theta_t )\\
    &A_{t+1}, \theta_{t+1} = \T(A_t, \theta_{t}') 
\end{aligned}
\end{equation}
in the hope that the final architecture $A_T$ and weights $\theta_T$ are a good approximation to the original objective~\eqref{eqn:ideal_obj}. The growth operator $\T$ is typically constrained to a neighborhood $\mathcal{N}(A_t)$ of the current architecture $A_t$, making local architecture modifications (\eg{}, adding neurons or layers).

In practice, the behavior of $\Opt$ heavily depends on the initialization $\theta_{t+1}$ of the transformed architecture, making it key to achieving good performance. The lottery ticket phenomenon~\citep{frankle_lottery_2019,chen_elastic_2021} highlights that the particular initialization and training path can matter as much as the final architecture, suggesting that growth methods leveraging a fixed set of initial weights have the potential to outperform NAS-like methods, which ignore this.

The abstract formulation of Equation~\eqref{eqn:grow_decomposition} leaves much of the growing problem unspecified, which is often summarized as \emph{where} to grow, \emph{when} to grow, \emph{how} to grow, and \emph{when to stop}. The focus in the literature has been overwhelmingly on \emph{how} to grow, which we term the \emph{neuron addition problem}: when adding new neurons, how many should we add, and how should we initialize them? 

\subsection{Motivations and applications}

Motivations for growing neural networks broadly fall into two categories. In the first, the end point is known: growth is a training strategy for reaching a predefined target architecture $A_T$. Growing is then motivated by faster training~\citep{evci_gradmax_2022,yuan_accelerated_2023}, by improved optimization dynamics~\citep{wu_when_2024,caillon_growing_2021}, or to account for changing data distributions such as in continual learning~\citep{yoon_lifelong_2018,li_learn_2019,yang_grown_2021}.

Alternatively, the endpoint architecture $A_T$ is unknown: growth is used as a frugal form of architecture search, where we would like to discover an architecture that is ``just large enough'' for the task at hand. Such \emph{frugal learning} is becoming increasingly important as achieving state-of-the-art performance increasingly relies on scaling model size and compute, with energy consumption and $\mathrm{CO_2}$ emissions increasing exponentially, outpacing improvements in hardware~\citep{thompson_deep_2021,morand_environmental_2025}. As a frugal method, growing is complementary to compression and pruning, which start from an oversized model and typically only remove parameters after training. While compression and pruning reduce inference costs, they still require training the larger model. The energy efficiency comparison with growing is nuanced, with recent work~\citep{boumendil_grow_2023} showing that the optimal strategy depends on the target accuracy and compute budget, with neither approach dominating across all scenarios.
    
Transformer growth is a prominent application of neural network growth, and is able to significantly reduce the training cost, by growing towards a target full-size architecture~\cite{gu_bert_2021,wang_ligo_2023}. In Sec.~\ref{sec:transformers}, we discuss the connection between transformer growth and the wider growing literature, and empirical evidence for their use.

This survey is \emph{the first modern overview of growing neural networks during training time}, the last one being close to thirty years ago~\citep{kwok_constructive_1997}. In contrast, other surveys target either sparsity and pruning in neural networks~\citep{sparsitySurvey}, inference-time dynamic architectures~\citep{dynamicSurvey}, or comparative studies restricted to transformers~\citep{pandey2024comparative}. Instead, this survey aims to unify the field of growing neural networks, scattered across various growth types, applications, and domains. In particular, we show that several key methods for neuron addition based on function geometry can be unified via a common optimization objective.

This survey is organized as follows. We formalize the neuron addition problem in Sec.~\ref{sec:neuron_addition_problem}, before describing the main classes of neuron initialization, i.e. \emph{how to grow}. We then discuss how the proposed criteria can be used to decide where to grow, when to grow, and when to stop in Sec.~\ref{sec:where_when}. We then broaden the scope beyond width and depth growth to discuss boosting (Sec.~\ref{sec:boosting}), sparse growth and grow-prune methods (Sec.~\ref{sec:sparse_prune_methods}), and then applications to continual learning and reinforcement learning in Sec.~\ref{sec:non_stationary}. Sec.~\ref{sec:results} aggregates a large body of width-growing and depth-growing literature to re-examine the degree to which growing satisfies its stated motivations, including how these conclusions transfer to transformer architecture growth. Finally, we discuss open questions and future directions.
